\section{Introduction}\label{sec:introduction}

Critical nearest-neighbor Ising spin interfaces have a conformally invariant
scaling limit:
with opposite boundary signs on two arcs, the limiting interface is chordal
\(\SLE_3\). We prove that this conclusion persists under every sufficiently
small finite-range perturbation of the contour energy that preserves the
symmetries of the square lattice. The perturbation may contain arbitrary
many-edge interactions and may have either sign. The boundary specification
is part of the result: interactions across the boundary remain present, and
the exterior may contain additional closed contours.

\subsection{The model and its boundary data}

Fix \(J>0\). Let \(\mathcal E^\bullet\) be the nearest-neighbor edges of
\(\Z^2\), and let \(\mathcal E\) be those of the dual lattice
\(\mathcal F=\Z^2+(1/2,1/2)\). Write \(e^\bullet\leftrightarrow e\) when the
primal and dual edges cross. A potential \(U\) assigns a real number to
each finite subset of \(\mathcal E\), with \(U(\varnothing)=0\).
We assume that \(U\) is invariant under lattice translations and the
rotations and reflections of the square lattice. Its range is finite:
there is \(R<\infty\) such that
\[
 U(S)=0
 \quad\text{if the Euclidean diameter of the union of the edges in \(S\)
 exceeds \(R\).}
\]
All values of \(U\) are fixed before the mesh is rescaled.

For a finite \(\Omega\subset\Z^2\), put
\[
 Q_\Omega=\bigcup_{v\in\Omega}\bigl(v+[-1/2,1/2]^2\bigr),
 \qquad D_\Omega=\operatorname{int}Q_\Omega.
\]
We require \(Q_\Omega\) to be the closure of a Jordan domain.
Let \(E_\Omega\subset\mathcal E\) consist of the edges crossing primal
bonds with at least one endpoint in \(\Omega\), and let \(F_\Omega\)
be their endpoints. Choose distinct
\(f_{\rm in},f_{\rm out}\in\partial D_\Omega\cap F_\Omega\).
For each mark choose a dual edge outside \(E_\Omega\) joining that mark
to a vertex outside \(Q_\Omega\); denote these edges by
\(e_{\rm in},e_{\rm out}\). They are the exterior stubs, and
\(\xi=((f_{\rm in},e_{\rm in}),(f_{\rm out},e_{\rm out}))\).

For a dual-edge set \(A\), let \(d_f(A)\) be its degree at \(f\).
The permitted interior contours are
\[
 \mathcal P_\Omega^\xi
 =
 \left\{P\subset E_\Omega:
 d_f(P)+\one_{\{f_{\rm in}\}}(f)+\one_{\{f_{\rm out}\}}(f)
 \text{ is even for every }f\in\mathcal F\right\}.
\]
We require this set to be nonempty. An admissible exterior contour is a
finite set \(P_o\subset\mathcal E\setminus E_\Omega\) containing the two
stubs and satisfying
\[
 d_f(P_o)+\one_{\{f_{\rm in}\}}(f)+\one_{\{f_{\rm out}\}}(f)
 \equiv0\pmod2
 \qquad(f\in\mathcal F).
\]
Thus \(P\cup P_o\) is even for every \(P\in\mathcal P_\Omega^\xi\).

For inverse temperature \(\beta>0\) and coupling \(\lambda\in\R\), the law is
\begin{equation}\label{eq:contour-law}
 \mu_{\Omega,\beta,\lambda}^{\xi,P_o}(P)
 =
 \frac{1}{Z_{\Omega,\beta,\lambda}^{\xi,P_o}}
 \exp\!\left\{
 -2\beta J\abs P
 -\beta\lambda
 \sum_{\varnothing\ne X\subseteq P}
 \sum_{Y\subseteq P_o} U(X\cup Y)
 \right\}.
\end{equation}
The subset sums include equality, and the second includes the empty set.
The normalizing constant is the sum over
\(\mathcal P_\Omega^\xi\). In particular, the boundary data retain every
interaction joining an occupied interior edge to occupied exterior edges.

The curve \(\gamma(P)\) uses \(P\) and the half of each exterior stub
between its mark and its midpoint. At degree two, join the two occupied
edges. At degree four, join north to east and south to west, also at a
mark if its half-stub raises the degree to four. Draw these connections
disjointly in disks of radius less than \(1/10\) around the dual vertices.
There is a unique component between the two stub midpoints; orient it
from the incoming midpoint to the outgoing midpoint. This is
\(\gamma(P)\). Its edges are each traversed at most once, and the resolved
curve is simple. Any fixed local drawing with these connections is allowed.

\subsection{The main theorem}

An oriented curve is a continuous map from \([0,1]\) into the plane,
considered modulo increasing reparametrization and zero distance for
\begin{equation}\label{eq:curve-metric}
 d_{\rm curv}(\gamma,\eta)
 =
 \inf_{\alpha,\rho}
 \sup_{t\in[0,1]}
 \abs{\gamma(\alpha(t))-\eta(\rho(t))},
\end{equation}
where \(\alpha,\rho\) range over increasing homeomorphisms of \([0,1]\).
The orientation distinguishes the two endpoints.

For a bounded Jordan domain \(D\) and distinct \(a,b\in\partial D\),
chordal \(\SLE_3\) from \(a\) to \(b\) is the conformal image of the
upper-half-plane Loewner trace with
\[
 \partial_t g_t(z)=\frac{2}{g_t(z)-\sqrt3 B_t},
 \qquad g_0(z)=z,
\]
where \(B\) is standard real Brownian motion; the conformal map sends
\(0\) to \(a\) and \(\infty\) to \(b\). This law is independent of the
choice of such map. The trace is simple and meets the Jordan boundary
only at its endpoints; see \citet{Schramm2000,RS2005}.

\begin{theorem}\label{thm:main}
Let \(J>0\), and let \(U\) satisfy the finite-range and symmetry assumptions
above. There is \(\lambda_0=\lambda_0(J,U)>0\) such that for every
\(\abs\lambda<\lambda_0\) there is a positive inverse temperature
\(\beta_c=\beta_c(J,U,\lambda)\), independent of domains and boundary data,
with
\[
 \beta_c(J,U,0)=\beta_0:=\frac{\log(1+\sqrt2)}{2J},
\]
for which the following holds.

Let \(D\) be any bounded Jordan domain, let \(a,b\) be distinct boundary
points, and let \(\delta_n\downarrow0\). For each \(n\), choose a domain
\(\Omega_n\), marks and stubs \(\xi_n\), and an admissible exterior
contour \(P_{o,n}\) as above. Suppose there are boundary homeomorphisms
\[
 \phi_n:\partial D\longrightarrow
       \partial(\delta_n D_{\Omega_n})
\]
such that
\[
 \sup_{z\in\partial D}\abs{\phi_n(z)-z}\longrightarrow0,\qquad
 \phi_n(a)=\delta_n f_{{\rm in},n},\qquad
 \phi_n(b)=\delta_n f_{{\rm out},n}.
\]
If \(P_n\) has law
\(\mu_{\Omega_n,\beta_c,\lambda}^{\xi_n,P_{o,n}}\), then
\(\delta_n\gamma(P_n)\) converges in distribution, in the metric
\eqref{eq:curve-metric}, to chordal \(\SLE_3\) in \(D\) from \(a\) to \(b\).
\end{theorem}

No smoothness of \(\partial D\) is assumed. The same unit-lattice
potential \(U\) is used for every \(n\), and the exterior contours may
vary arbitrarily with \(n\). The theorem concerns the full curve law;
it does not assume convergence of an observable or a tightness
conjecture for the perturbed model.

\subsection{Historical context}

Schramm's Loewner description identifies conformally invariant random
interfaces through a one-dimensional driving process
\citep{Schramm2000}. Smirnov proved conformal covariance of the critical
FK-Ising fermionic observable on the square lattice
\citep[Theorem~2.2]{Smirnov2010}; Chelkak and Smirnov developed the spin
and FK observables on isoradial graphs
\citep[Theorems~A--B]{ChelkakSmirnov2012}. Together with geometric
crossing control, these observables lead to convergence of the critical
nearest-neighbor spin interface to \(\SLE_3\)
\citep[Theorem~1]{CDHKS2014}. Conformal covariance of spin correlations
and convergence of nested loop collections to \(\CLE_3\) are established
in \citet{CHI2015,BH2019}. The spin alternating-arm exponents,
including the four-arm value \(21/8\), are given by \citet{Wu2018}.

For non-integrable perturbations, constructive renormalization has
established universality for several correlation observables.
Giuliani, Greenblatt, and Mastropietro prove full-plane energy-correlation
scaling limits for weak, rotation-invariant finite-range pair
perturbations \citep{GGM2012}. Antinucci, Giuliani, and Greenblatt
extend energy universality to finite cylinders with periodic boundary
conditions in one direction and free boundary conditions in the other,
allowing translation-invariant finite-range even multispin interactions
\citep[Theorem~1.1 and Corollary~1.2]{AGG2023Cylinder}.
Cava, Giuliani, and Greenblatt prove boundary-spin scaling limits in
the half-plane with free boundary conditions, including the
renormalized Pfaffian limit under weak translation-invariant finite-range
even perturbations
\citep[Theorem~1.1]{CGG2025Boundary}.
These results treat either sign of a sufficiently small perturbation
and give precise information in the stated plane, cylinder, and half-plane
geometries.

A different approach uses random currents. For ferromagnetic,
translation-invariant, square-symmetric finite-range pair interactions
whose interaction graph is connected,
\citet[Theorem~1.2]{AizenmanDuminilCopinTassionWarzel2019}
prove an asymptotic Pfaffian structure of critical boundary-spin
correlations in the half-plane, without a small-perturbation assumption.
Its relative Pfaffian relation and the preceding correlation scaling
limits concern observables and hypotheses different from the full
signed contour-potential interface law considered here.

A close predecessor for the interface problem is the work of
\citet{GreenblattPeltola2024} on non-integrable finite-range
interactions. Their Theorem~4.4 expresses an exploration local martingale as a ratio
of Grassmann correlations, and their Conjecture~4.5 formulates fermionic
asymptotics,
local renormalization, and boundary stability that would yield an
\(\SLE_3\) limit. Their Section~5 explains why estimates in fixed
geometries must be extended to the irregular boundaries created by
exploration. We establish the interface limit for the symmetric class
in Theorem~\ref{thm:main} through bounded spin observables, sparse
changes, and restricted boundary comparisons. The specific fermionic
asymptotics in their conjecture are not a conclusion of this argument.

There is a complementary geometric lineage. Multi-crossing estimates
control random-curve regularity in
\citet[Hypothesis~H1 and Theorem~1.1]{AizenmanBurchard1999}, and
conditional bounds on unforced crossings yield tightness and Loewner limits in
\citet[Condition~G2 and Theorem~1.5]{KemppainenSmirnov2017}.
The distinction between driving-function
convergence and strong curve convergence is developed in
\citet[Theorem~1.2 and Example~2.1]{SheffieldSun2012}.
Our geometric input is a mesh-first interior
four-passage estimate. We prove the trace-to-curve implication needed
here directly by excluding collapsing reverse traversals.

\subsection{Proof strategy}

We work with spins by assigning plus at infinity and changing sign each
time a primal path crosses an edge of the even contour set \(P\cup P_o\).
This assigns a spin \(\sigma_v\in\{-1,+1\}\) to each lattice vertex;
the spins outside \(\Omega\) are fixed by \(P_o\). A prescribed spin
is called a \emph{pin}. Expanding each
contour interaction in these spins gives a finite-range even spin
interaction. Section~\ref{subsec:spin-formulation} proves the exact
identity of the two conditional laws, including all boundary interactions.
Our first comparison is between the periodic version of this bulk
interaction and the critical nearest-neighbor law on a large periodic
square.

Section~\ref{sec:pure} specifies the pure-model estimates used in that
comparison. Besides the public results above, we use the buffered
likelihood comparisons, ordered-source transmission, nearby-line
prediction, and complete-resolution results of the companion article
\emph{Buffered comparison and stopping-band resolution in critical
Ising}~\cite{OpenAIStatic2026}. These are equilibrium statements about
the ordinary critical model. Their hypotheses distinguish arbitrary
common pins in likelihood comparisons from the free neighborhoods
needed to transmit a magnetic response.

The linear transfer takes a bounded spin function supported in a square,
conditions it under the pure law on the boundary spins of a concentric square
larger by a fixed ratio \(L\), and multiplies by \(L^2\) to normalize
an interaction density. The first obstacle is that repeated transfer
creates functions of more and more microscopic spins. Formulas for a fixed number of field
insertions do not control all such functions. Camia, Garban, and Newman
construct smeared magnetization-field limits from spin correlations
\citep[Section~3]{CGN2015}; Section~\ref{sec:transfer} proves the further
uniform approximation needed here. We observe each spin through a weak
Gaussian signal. Conditional on these observations the pure law remains
ferromagnetic. Ordered-source transmission and prediction from nearby
spin sums show that the remaining uncertainty has little effect across
a positive gap. Only finitely many spatial averages of the observation
are then needed. This reduces the transfer of arbitrary bounded spin
functions to the explicit continuum spin formulas. Reflection positivity
and a finite determinant identity show that, among even square-invariant
interaction densities, exactly one scalar direction can grow under this
transfer.

Section~\ref{sec:flow} turns the linear transfer into an exact flow of
local interactions. A finite-state identity explains the construction:
moving a potential toward its conditional expectation can be reversed
by random local refreshes at rates independent of the spins. We control
the interactions produced at successive scales. Positive long-distance
energy covariance shows that changing temperature moves transversely
to the set of initial interactions whose trajectories decay. We can
therefore choose the inverse temperature to obtain such a trajectory.
This choice uses only bulk data and is
therefore independent of the domain. Reversing the flow couples the
resulting torus spin law to the pure law by changes inside independently
marked squares. The marking probabilities decay with scale; the spins
assigned inside marked squares remain unrestricted.

Section~\ref{sec:stability} proves that these possible changes preserve
every fixed finite collection of polygonal crossing, circuit, and
four-passage tests with probability tending to one. The geometric reason
has three forms. An edit in the interior of a test can change its answer
only if four long contour portions reach the edit. Near a straight test
edge two portions must remain on one side, and near a corner one portion
must extend to a fixed distance. The corresponding probabilities decay
faster than the numbers of possible locations grow. The strict bulk
threshold comes from the spin four-arm exponent \(21/8>2\); the
half-plane estimate is derived here. A bootstrap over the edit scales
makes the conclusion valid for every assignment in all marked squares,
including assignments chosen after examining the pure configuration.

To pass from periodic squares to the prescribed domain, we need comparison
under boundary conditions. Signed perturbations need not preserve
ferromagnetic comparison for the physical spins. Section~\ref{sec:comparison}
represents each small interaction by additional binary variables on
nearest-neighbor bonds. Their local conditional probabilities have strict
margins under specified changes from minus pins to free spins or from free
spins to plus pins. These margins provide the boundary comparisons used
here. Combined with the bulk crossing tests, they also bound the
pointwise density ratio between a target-domain spin marginal on a
square and the perturbed torus marginal on that square, when a fixed
free buffer separates the square from all pins.

Finally, Section~\ref{sec:curves} constructs two displaced pure comparison
domains, favoring the opposite colors. Finite collections of colored strip
crossings transfer barriers near their interfaces to the target. In any
subsequential limit, the barriers and connectedness of the target trace
order the two limiting simple chords. Both chords have the same
\(\SLE_3\) law, so they coincide and identify the trace. A persistent backward traversal near a
simple chord would create six passages through a small interior annulus.
The four-passage estimate excludes this event and yields convergence in
the full oriented curve metric~\eqref{eq:curve-metric}.

The inverse temperature is selected by the decaying trajectories of the
conditional flow in Section~\ref{sec:flow}. Its local uniqueness refers to
that construction. No identification with a temperature branch selected
by a different scale map, or joint field-and-interface scaling limit,
is required for Theorem~\ref{thm:main}.

\subsection{Exact spin formulation}\label{subsec:spin-formulation}


The proof uses spin conditional expectations in the bulk. We record the
exact passage from \eqref{eq:contour-law} to spins, including the exterior
data. The contour-potential formulation with exterior interactions is
also used by \citet[Section~4]{GreenblattPeltola2024}. For a primal bond crossing \(e\), write
\[
 d_e(\sigma)=\frac{1-\sigma_v\sigma_w}{2},
 \qquad e^\bullet=\{v,w\},\quad \sigma\in\{-1,1\}^{\Z^2}.
\]
The formal spin Hamiltonian is
\begin{equation}\label{eq:intro-spin-hamiltonian}
 \mathcal H_{\beta,\lambda}(\sigma)
 =
 -\beta J\sum_{\{v,w\}\in\mathcal E^\bullet}\sigma_v\sigma_w
 +\beta\lambda
  \sum_{\varnothing\ne S\subset\mathcal E}
        U(S)\prod_{e\in S}d_e(\sigma).
\end{equation}
Only terms depending on the spins in the chosen finite domain enter its
conditional law. All resulting sums are finite.

\begin{lemma}\label{lem:contour-spin}
For the boundary data above, every even completion \(P\cup P_o\)
determines a unique spin configuration that is plus sufficiently far
away and has disagreement set \(P\cup P_o\). Its restriction
\(\tau\) to \(\Z^2\setminus\Omega\) is independent of \(P\).
The map \(P\mapsto\sigma|_\Omega\) is a bijection between
\(\mathcal P_\Omega^\xi\) and all spin assignments in \(\Omega\).
Under this bijection, \eqref{eq:contour-law} is precisely the Gibbs
conditional law for \eqref{eq:intro-spin-hamiltonian} with exterior
spins \(\tau\). Denote this law by
\(\nu_{\Omega,\beta,\lambda}^{\tau}\).

The exterior nearest neighbors along each of the two boundary arcs
between the marks have a constant sign, and the signs of the two arcs
are opposite. The interaction in
\eqref{eq:intro-spin-hamiltonian} is finite-range, invariant under global
spin flip and the square symmetries, and has finitely many local
interaction types up to translation.
\end{lemma}

\begin{proof}
A finite even dual subgraph has even intersection parity with every
closed primal walk. One way to see this is to reduce a closed walk
modulo two to elementary primal faces; the intersection parity around
such a face is the degree of the dual vertex it encloses. Starting with
a plus spin at infinity, flip signs on crossing an occupied dual edge.
Path independence follows from this parity statement. Finiteness fixes
the plus sign in the unbounded region and gives uniqueness.

The exterior cell centers \(\Z^2\setminus\Omega\) form a connected
nearest-neighbor graph. Indeed the complement of a Jordan polyomino has
no bounded component, and a path in its exterior can be perturbed to
cross cell sides away from vertices, giving a path of exterior cells.
The Jordan condition excludes a diagonal pinch. A primal path using
only exterior vertices crosses only dual edges outside \(E_\Omega\).
Their occupied set is \(P_o\), so all exterior signs are fixed
independently of \(P\).

Conversely, extend any assignment in \(\Omega\) by these fixed exterior
spins. Its disagreement set outside \(E_\Omega\) is \(P_o\), and its
complete dual disagreement graph is even. The interior part therefore
belongs to \(\mathcal P_\Omega^\xi\). These two constructions are inverse.

For the signed-arc assertion, traverse the polygonal boundary. At a
flat boundary vertex, the two successive exterior neighboring cells
are separated by one exterior primal bond. At an unmarked vertex the
crossing exterior dual edge is absent by the parity condition, so the
two exterior signs agree. At a convex corner, the corresponding
exterior path has two bonds; its crossing parity is again even at
an unmarked vertex. At a reflex corner, the two boundary edges have
the same adjacent exterior cell. At either mark the exterior parity
is odd, and the sign changes. Reflex corners cannot support an exterior
stub of the required kind. Thus there are exactly the two asserted
constant, opposite arc signs.

Now fix an interior spin configuration and its contour \(P\).
A nonconstant interaction term in
\eqref{eq:intro-spin-hamiltonian} has a support \(S\) meeting
\(E_\Omega\). Write
\[
 X=S\cap E_\Omega,\qquad Y=S\setminus E_\Omega.
\]
The product of disagreement indicators is one exactly when
\(X\subseteq P\) and \(Y\subseteq P_o\). Summing over all such supports
gives precisely the double sum in \eqref{eq:contour-law}.
Terms supported entirely outside \(E_\Omega\) are constant in the
interior spins. The nearest-neighbor energy is \(2\beta J\abs P\)
up to an exterior-dependent constant. This proves equality of the
conditional laws.

Finally, choose an edge in a nonempty support of \(U\). Finite range
leaves only finitely many possible other edges relative to it, hence
only finitely many support types up to translation. Expanding the
products of \(d_e\) gives finite local spin polynomials, invariant
under spin flip and under the prescribed square symmetries.
\end{proof}

In particular, exterior contours do not give an additional bulk
parameter. Their only effect is through the actual exterior spins
in the conditional law. This observation does not permit deleting
interactions across a narrow fjord or replacing distant exterior
spins by their adjacent arc signs.

